\section{Economía como ciencia y aplicación al comercio internacional}
\label{sec:economia-ciencia-comercio}

La economía, como ciencia social, trabaja con modelos y supuestos para explicar comportamientos complejos. Los modelos no pretenden replicar la realidad completa; buscan aislar el efecto de una variable sobre otra para que el análisis sea más claro.

\subsection{Modelos y supuestos}

\begin{concepto}{Modelo económico}
Un modelo es una representación simplificada de la realidad que permite estudiar una relación causal o un fenómeno concreto.
\end{concepto}

Los economistas usan supuestos para reducir la complejidad. Por ejemplo, al analizar la demanda, se supone que todos los demás factores se mantienen constantes, salvo el precio del bien que se estudia.

\begin{ejemplo}
Si se analiza la relación entre el precio del pan y la cantidad demandada, se mantiene constante el ingreso de las personas, el precio de la mantequilla y las preferencias del consumidor. Así se puede estudiar el efecto del precio en aislado.
\end{ejemplo}

\subsection{Ceteris paribus}

\begin{concepto}{Ceteris paribus}
Expresión latina que significa “manteniendo todo lo demás constante”.
\end{concepto}

Es una herramienta fundamental en economía porque permite comparar situaciones de manera ordenada. Sin este supuesto, el análisis sería confuso, porque todos los factores cambiarían a la vez.

\subsection{Correlación versus causalidad}

\begin{itemize}
    \item \textbf{Correlación:} dos variables se mueven juntas.
    \item \textbf{Causalidad:} un cambio en una variable provoca un cambio en la otra.
\end{itemize}

\begin{advertencia}{Error clásico}
Que dos fenómenos se muevan juntos no significa que uno cause al otro. A veces hay una variable omitida que explica la relación.
\end{advertencia}

\begin{ejemplo}
Se observa que las ventas de helados y los ahogamientos aumentan en verano. No significa que los helados causen ahogamientos; lo que cambia es la temperatura del clima.
\end{ejemplo}

\subsection{Economía positiva y normativa}

\begin{itemize}
    \item \textbf{Economía positiva:} describe cómo funciona el mundo. Busca hechos y relaciones verificables.
    \item \textbf{Economía normativa:} evalúa cómo debería ser el mundo, según juicios de valor.
\end{itemize}

\begin{ejemplo}
- Positiva: “Si aumenta el precio del combustible, la cantidad demandada cae.”
- Normativa: “El Estado debería subsidiar el transporte para reducir la contaminación.”
\end{ejemplo}

\subsection{La Frontera de Posibilidades de Producción (FPP)}

La FPP muestra todas las combinaciones posibles de producción de dos bienes que una economía puede producir, usando eficientemente sus recursos y tecnología disponibles.

\begin{concepto}{FPP}
La Frontera de Posibilidades de Producción representa las opciones máximas de producción de una economía.
\end{concepto}

\subsubsection{Supuestos del modelo FPP}

Para construir la FPP, se asume que:
\begin{enumerate}
    \item la economía produce solo dos bienes,
    \item tiene recursos fijos,
    \item usa todos los recursos de manera eficiente,
    \item la tecnología es dada.
\end{enumerate}

\begin{itemize}
    \item \textbf{Puntos sobre la frontera:} producción eficiente.
    \item \textbf{Puntos dentro de la frontera:} producción ineficiente.
    \item \textbf{Puntos fuera de la frontera:} inalcanzables con la tecnología y recursos actuales.
\end{itemize}

\begin{ejemplo}
Supongamos una economía que produce solo alimentos y ropa. Si usa todos sus recursos eficientemente, puede producir más alimentos o más ropa, pero no ambas cantidades máximas al mismo tiempo. La FPP representa ese equilibrio entre usos alternativos.
\end{ejemplo}

\subsection{Costo de oportunidad en la FPP}

La pendiente de la FPP mide el costo de oportunidad. Es decir, cuántas unidades de un bien se deben sacrificar para producir una unidad más del otro bien.

\begin{concepto}{Costo de oportunidad en la FPP}
Es la cantidad de un bien que se deja de producir para obtener más del otro.
\end{concepto}

\begin{ejemplo}
Si una economía pasa de producir 100 unidades de alimentos y 80 de ropa a producir 110 alimentos y 60 ropa, el costo de oportunidad de producir 10 alimentos adicionales fue 20 unidades de ropa.
\end{ejemplo}

\subsection{Formas de la FPP}

\begin{itemize}
    \item \textbf{FPP recta:} implica rendimientos constantes.
    \item \textbf{FPP cóncava:} implica rendimientos decrecientes.
\end{itemize}

\begin{concepto}{Rendimientos decrecientes}
A medida que se especializa más en un bien, cada unidad adicional requiere renunciar a una mayor cantidad del otro bien.
\end{concepto}

\begin{ejemplo}
Si un país está produciendo mucho vino y poco trigo, y decide producir un poco más de trigo, tendrá que sacrificar una gran cantidad de vino, porque los recursos menos adecuados para trigo ya están siendo usados en vino.
\end{ejemplo}

\subsection{Desplazamientos de la FPP}

La FPP se desplaza hacia afuera cuando hay crecimiento económico, por ejemplo:
\begin{itemize}
    \item más trabajadores,
    \item más capital,
    \item progreso tecnológico,
    \item mayor productividad.
\end{itemize}

\begin{observacion}{Importante}
Un cambio en preferencias o precios no desplaza la FPP. Solo cambia el punto de producción elegido dentro o sobre la frontera.
\end{observacion}

\subsection{Comercio internacional y ventajas}

El comercio internacional se explica a partir de las ventajas comparativas. El intercambio permite que cada país produzca aquello en lo que es relativamente mejor o más eficiente.

\subsubsection{Ventaja absoluta}

\begin{concepto}{Ventaja absoluta}
Un país tiene ventaja absoluta si puede producir un bien usando menos recursos que otro país.
\end{concepto}

\begin{ejemplo}
Si Chile puede producir 1 tonelada de vino con menos trabajo que Argentina, Chile tiene ventaja absoluta en vino.
\end{ejemplo}

\subsubsection{Ventaja comparativa}

\begin{concepto}{Ventaja comparativa}
Un país tiene ventaja comparativa si puede producir un bien con un menor costo de oportunidad.
\end{concepto}

La ventaja comparativa es la base del comercio. Incluso si un país es mejor en todo, puede beneficiarse del comercio especializándose en lo que hace relativamente mejor.

\begin{ejemplo}
Supongamos que dos países pueden producir trigo y vino. Si Chile tiene menor costo de oportunidad en vino y Argentina en trigo, ambos ganan si Chile se especializa en vino y Argentina en trigo y luego comercian.
\end{ejemplo}

\subsection{Autarquía y comercio}

\begin{itemize}
    \item \textbf{Autarquía:} economía cerrada, sin comercio exterior.
    \item \textbf{Comercio:} intercambio entre países, basado en especialización.
\end{itemize}

\begin{concepto}{Ganancias del comercio}
El comercio aumenta el bienestar porque permite especializarse, producir más eficientemente y consumir más variedad.
\end{concepto}

\begin{ejemplo}
Si un país puede producir trigo muy eficientemente y otro vino muy eficientemente, ambos mejoran si cada uno se especializa y luego intercambia bienes.
\end{ejemplo}

\subsection{Rango de precios relativos}

Para que exista comercio mutuamente beneficioso, el precio de intercambio debe ubicarse entre los costos de oportunidad de ambos países.

\begin{ejemplo}
Si el costo de oportunidad de Chile de producir 1 vino es 2 trigo y el de Argentina es 3 trigo, un comercio aceptable será a un precio intermedio, por ejemplo 2.5 trigo por 1 vino.
\end{ejemplo}

\subsection{Frontera de Posibilidades de Consumo (FPC)}

Cuando los países comercian, ya no están limitados únicamente a su producción interna. Su nueva restricción es la FPC, determinada por el precio internacional.

\begin{concepto}{FPC}
La FPC muestra lo que puede consumir un país cuando se especializa y comercia.
\end{concepto}

\begin{ejemplo}
Si Chile se especializa en vino y exporta parte de su producción, puede consumir más trigo del que podría producir por sí solo. Eso representa un aumento del bienestar.
\end{ejemplo}

\subsection{Ejercicio resuelto: ventaja comparativa}

\begin{ejercicio}{Robinson y Viernes}
Robinson y Viernes trabajan 40 horas a la semana y producen cocos y langostas.

- Robinson tarda 2 horas en un coco y 8 horas en una langosta.
- Viernes tarda 30 horas en un coco y 10 horas en una langosta.

\textbf{Paso 1:} cálculo del costo de oportunidad.

Robinson:
- 1 coco cuesta 0.25 langostas.
- 1 langosta cuesta 4 cocos.

Viernes:
- 1 coco cuesta 3 langostas.
- 1 langosta cuesta 0.33 cocos.

\textbf{Paso 2:} quien tiene ventaja comparativa.

- Robinson tiene menor costo de oportunidad en cocos.
- Viernes tiene menor costo de oportunidad en langostas.

\textbf{Conclusión:}
Robinson debe especializarse en cocos y Viernes en langostas. Luego comercian y ambos salen mejor.
\end{ejercicio}

\subsection{Resumen final del bloque}

La economía funciona como una ciencia social basada en modelos, supuestos y análisis causal. La Frontera de Posibilidades de Producción ayuda a visualizar la escasez y el costo de oportunidad. A partir de allí, el comercio internacional explica por qué los países se benefician al especializarse en aquello que producen relativamente mejor.

\begin{resumen}{Lo clave del bloque}
- Los modelos simplifican la realidad.
- Ceteris paribus permite analizar un cambio aislado.
- La FPP muestra el costo de oportunidad.
- El comercio surge de la ventaja comparativa.
- La especialización internacional aumenta bienestar.
\end{resumen}
